% !TeX root = ../../Book2.tex
\section{字符与对偶群}
\label{b2:sec:I203}

在实直线上，指数函数满足
\[
 e^{2\pi\mathrm i\xi(x+y)}=e^{2\pi\mathrm i\xi x}e^{2\pi\mathrm i\xi y}.
\]
这个等式使平移在 Fourier 变换后成为乘法。周期函数的指数基、离散 Fourier 变换中的单位根也满足同一关系。把这种关系作为定义，就能在一个交换群上寻找它自己的“频率”。频率所组成的空间可能离散，也可能连续；它的拓扑会决定 Fourier 展开应当写成求和还是积分。

\subsection{连续字符与字符分离}
\label{b2:subsec:I20301}

\begin{definition}[连续字符]
\label{b2:def:I3-character}
设 \(G\) 为第二可数局部紧 Hausdorff 交换群，群运算记作加法。记
\[
 \mathbb S^1=\{z\in\mathbb C:|z|=1\},\qquad
 \mathbb T=\mathbb R/\mathbb Z.
\]
从 \(G\) 到乘法群 \(\mathbb S^1\) 的连续同态称为 \(G\) 的\term[zifu]{字符}[character]，即连续函数 \(\chi:G\to\mathbb S^1\) 满足
\[
 \chi(x+y)=\chi(x)\chi(y)\qquad(x,y\in G).
\]
所有字符组成的集合记为 \(\widehat G\)。恒等于 \(1\) 的字符称为平凡字符，记为 \(\mathbf1\)。
\end{definition}
\symidx{\(\widehat G\)：局部紧交换群的对偶群}

这里 \(\mathbb T\) 使用加法，\(\mathbb S^1\) 使用乘法；映射 \(t+\mathbb Z\mapsto e^{2\pi\mathrm it}\) 将二者识别。保留两种写法，可以同时看清周期变量和复数相位。由同态关系可得
\[
 \chi(0)=1,\qquad \chi(-x)=\chi(x)^{-1}=\overline{\chi(x)}.
\]
字符的逐点乘积仍为字符，逐点共轭给出逆元。因此 \(\widehat G\) 本身是交换群。对平移 \(L_af(x)=f(x-a)\)，有
\[
 L_a\chi=\overline{\chi(a)}\chi.
\]
这说明字符是全部平移算子的共同特征函数；频率就是这些特征值随 \(a\) 变化的记录。

先求出最熟悉的三个群的全部字符。实直线的情形有一个值得展开的步骤：连续的复相位可以在单位元附近选定辐角，随后同态关系把这个局部选择延伸到全直线。

\begin{proposition}[实直线的字符]
\label{b2:prop:I3-real-characters}
实加法群 \(\mathbb R\) 的每个连续字符都唯一地写成
\[
 \chi_\xi(x)=e^{2\pi\mathrm i\xi x}\qquad(x\in\mathbb R)
\]
的形式，其中 \(\xi\in\mathbb R\)。
\end{proposition}
\begin{proof}
设 \(\chi\) 是这样的字符。由 \(\chi(0)=1\) 和连续性，可选 \(\delta>0\)，使 \(|x|<\delta\) 时 \(\chi(x)\) 落在弧
\(\{e^{2\pi\mathrm it}:|t|<1/8\}\) 内。指数映射在该区间上有连续逆，故存在连续函数
\[
 a:(-\delta,\delta)\longrightarrow(-1/8,1/8),\qquad
 a(0)=0,\quad \chi(x)=e^{2\pi\mathrm ia(x)}.
\]
若 \(x,y,x+y\) 都在此区间，同态关系表明
\(a(x+y)-a(x)-a(y)\) 是整数；它的绝对值小于 \(3/8\)，所以只能为零。于是 \(a\) 在所有可作加法的局部位置满足可加性。

对任意 \(t\in\mathbb R\)，选正整数 \(m\) 使 \(|t|/m<\delta\)，定义
\[
 A(t)=m\,a(t/m).
\]
这个值与 \(m\) 无关。事实上，如果 \(n\) 也是可选整数，沿着从 \(0\) 到 \(t/m\) 的 \(n\) 个小步反复使用局部可加性，得到
\[
 a(t/m)=n\,a(t/(mn)).
\]
所有中间点的绝对值都不超过 \(|t|/m<\delta\)，所以每一步都合法。同样
\(a(t/n)=m\,a(t/(mn))\)，两式给出所需独立性。

给定 \(s,t\)，取同一个充分大的 \(m\)，使 \(s/m,t/m,(s+t)/m\) 都在上述区间。由局部可加性，
\[
 A(s+t)=A(s)+A(t).
\]
在 \((-\delta,\delta)\) 上可取 \(m=1\)，故 \(A=a\) 在零点连续。连续可加函数必为线性函数：令 \(\xi=A(1)\)，先由整数倍和整数除法得到 \(A(q)=q\xi\) 对全部有理数 \(q\) 成立，再用有理数逼近和连续性得到 \(A(t)=\xi t\)。最后，
\[
 \chi(t)=\chi(t/m)^m=e^{2\pi\mathrm iA(t)}=e^{2\pi\mathrm i\xi t}.
\]
若 \(\chi_\xi=\chi_\eta\) 且 \(\xi\ne\eta\)，取 \(x=1/(2(\xi-\eta))\) 就会得到 \(e^{2\pi\mathrm i(\xi-\eta)x}=-1\)，矛盾。这也证明了参数的唯一性。
\end{proof}

\begin{proposition}[整数、圆周和循环群的字符]
\label{b2:prop:I3-model-duals}
赋予 \(\mathbb Z\) 离散拓扑、\(\mathbb T=\mathbb R/\mathbb Z\) 商拓扑，并设 \(N\ge1\) 为整数。下列公式分别给出各群的全部连续字符，而且所列参数唯一：
\begin{align*}
 G=\mathbb Z:&\quad \chi_\theta(n)=e^{2\pi\mathrm in\theta},
     &&\theta\in\mathbb T;\\
 G=\mathbb T:&\quad \chi_n(t)=e^{2\pi\mathrm int},
     &&n\in\mathbb Z;\\
 G=\mathbb Z/N\mathbb Z:&\quad \chi_k(j)=e^{2\pi\mathrm ikj/N},
     &&k\in\mathbb Z/N\mathbb Z.
\end{align*}
此外，若 \(G_1,\ldots,G_d\) 是第二可数局部紧 Hausdorff 交换群，则 \(G_1\times\cdots\times G_d\) 的每个字符都唯一地写成
\[
 \chi(x_1,\ldots,x_d)=\prod_{j=1}^d\chi_j(x_j),\qquad
 \chi_j\in\widehat G_j.
\]
\end{proposition}
\begin{proof}
\(\mathbb Z\) 上的同态完全由 \(\chi(1)\) 决定：\(\chi(n)=\chi(1)^n\)。每个单位复数唯一对应一个 \(\theta\in\mathbb T\)，而离散域上的映射自动连续。这得到第一项。

设 \(\chi\) 是 \(\mathbb T\) 上的字符，\(q:\mathbb R\to\mathbb T\) 是商映射。由前一命题，\(\chi\circ q(x)=e^{2\pi\mathrm i\xi x}\)。因为 \(q(x+1)=q(x)\)，必须有 \(e^{2\pi\mathrm i\xi}=1\)，即 \(\xi=n\in\mathbb Z\)。反过来，整数 \(n\) 保证公式不随代表元 \(t\) 改变，并由商拓扑得到连续性。

循环群上同样只需指定生成元的值。关系 \(N\cdot1=0\) 要求 \(\chi(1)^N=1\)，所以其值恰为 \(N\) 次单位根，给出第三项。

最后，将乘积群中的元素分解为各坐标轴元素之和。对字符使用同态关系，得到所列乘积公式，其中 \(\chi_j\) 是 \(\chi\) 在第 \(j\) 个坐标轴上的限制。各坐标嵌入连续，故限制是连续字符；反方向的乘积也连续。限制到坐标轴还给出唯一性。
\end{proof}

例如，\(\mathbb R^d\) 上的字符恰为 \(e^{2\pi\mathrm ix\cdot\xi}\)，\(\mathbb T^d\) 上的字符恰为 \(e^{2\pi\mathrm in\cdot t}\)，\(\mathbb Z^d\) 上的字符恰为 \(e^{2\pi\mathrm in\cdot\theta}\)。对一个已写成循环群乘积的有限群
\[
 F=\prod_{j=1}^r\mathbb Z/N_j\mathbb Z,
\]
相应字符为
\[
 \chi_k(x)=\exp\left(2\pi\mathrm i\sum_{j=1}^r\frac{k_jx_j}{N_j}\right),
 \qquad k\in\prod_{j=1}^r\mathbb Z/N_j\mathbb Z.
\]
它们把离散 Fourier 变换中的相位因子组织为一个群。

字符\term[zifufenli]{分离点}[separate points]，是说对任意 \(x\ne y\)，存在 \(\chi\) 使 \(\chi(x)\ne\chi(y)\)。由于 \(\chi(x)\overline{\chi(y)}=\chi(x-y)\)，只需对每个非零 \(z\) 找到 \(\chi(z)\ne1\)。上述模型中可以直接完成这一步。在 \(\mathbb R^d\) 上，若 \(z_j\ne0\)，取 \(\xi_j=1/(2z_j)\)、其余坐标为零，字符值便为 \(-1\)；在 \(\mathbb Z^d\) 上作同样的选择并把参数取模 \(1\)。在 \(\mathbb T^d\) 上，若 \(z_j\ne0\) 于 \(\mathbb T\)，取频率向量为第 \(j\) 个标准基向量。在循环群乘积中也只需选一个非零坐标及该坐标的基本字符。一般局部紧交换群的字符分离是 Pontryagin 对偶理论的一部分，将在\secref{b2:sec:I205}中以完整条件陈述。

\subsection{对偶群的拓扑与基本例子}
\label{b2:subsec:I20302}

频率接近，应当使它们在每个固定的有限观测区域内产生接近的相位。例如 \(\xi\to\eta\) 时，\(e^{2\pi\mathrm ix\xi}\) 在有界 \(x\) 区间上一致趋于 \(e^{2\pi\mathrm ix\eta}\)；在整条实线上取上确界却通常仍得到 \(2\)。因此，对偶群采用紧集上的一致收敛。

\begin{definition}[对偶群的紧开拓扑]
\label{b2:def:I3-dual-topology}
设 \(G\) 为第二可数局部紧 Hausdorff 交换群。对 \(\chi\in\widehat G\)、紧集 \(K\subset G\) 和 \(\varepsilon>0\)，令
\[
 U(\chi;K,\varepsilon)=
 \left\{\psi\in\widehat G:
   \sup_{x\in K}|\psi(x)-\chi(x)|<\varepsilon\right\}.
\]
由这些集合构成的邻域基定义 \(\widehat G\) 的\term[jinkaituopu]{紧开拓扑}[compact-open topology]；赋予此拓扑的字符群称为 \(G\) 的\term[duiouqun]{对偶群}[dual group]。
\end{definition}

有限个紧集的并仍然紧，故把若干 \(U(\chi;K_j,\varepsilon_j)\) 同时作为条件时，可以用 \(K=\bigcup_jK_j\) 和 \(\varepsilon=\min_j\varepsilon_j\) 缩小邻域。这个具体描述足以直接检验下面的连续性。

\begin{proposition}[对偶群的群拓扑]
\label{b2:prop:I3-dual-topology}
设 \(G\) 为第二可数局部紧 Hausdorff 交换群。紧开拓扑使 \(\widehat G\) 成为 Hausdorff 交换拓扑群，且对每个固定的 \(x\in G\)，求值映射 \(\chi\mapsto\chi(x)\) 连续。
\end{proposition}
\begin{proof}
两个不同字符在某个 \(x\) 处取值不同。取紧集 \(K=\{x\}\)，再取小于两值距离一半的邻域半径，就能用互不相交的邻域分开它们。相同的单点紧集说明求值连续。

对单位模复数，三角不等式给出
\[
 |\chi_1(x)\psi_1(x)-\chi_0(x)\psi_0(x)|
 \le |\chi_1(x)-\chi_0(x)|+|\psi_1(x)-\psi_0(x)|.
\]
在任意紧集上取上确界，便知若两个因子各改变小于 \(\varepsilon/2\)，乘积改变小于 \(\varepsilon\)。故群乘法连续。逆映射是逐点共轭，且
\(|\overline{\chi_1(x)}-\overline{\chi_0(x)}|=|\chi_1(x)-\chi_0(x)|\)，所以也连续。
\end{proof}

\begin{proposition}[基本对偶的拓扑识别]
\label{b2:prop:I3-model-topology}
把 \(\mathbb Z\) 和有限循环群赋予离散拓扑，把 \(\mathbb R\) 赋予通常拓扑，把 \(\mathbb T\) 赋予商拓扑。\cref{b2:prop:I3-real-characters}和\cref{b2:prop:I3-model-duals}给出的对应是拓扑群同构：
\[
 \widehat{\mathbb R\kern0pt}\simeq\mathbb R,\qquad
 \widehat{\mathbb Z\kern0pt}\simeq\mathbb T,\qquad
 \widehat{\mathbb T\kern0pt}\simeq\mathbb Z,\qquad
 \widehat{\mathbb Z/N\mathbb Z\kern0pt}\simeq\mathbb Z/N\mathbb Z.
\]
对有限个第二可数局部紧 Hausdorff 交换群，还成立拓扑群同构
\[
 \widehat{G_1\times\cdots\times G_d}
 \simeq\widehat G_1\times\cdots\times\widehat G_d.
\]
\end{proposition}
\begin{proof}
在 \(\mathbb R\) 上，紧集 \(K\) 包含在某个 \([-M,M]\) 内，故
\[
 \sup_{x\in K}|e^{2\pi\mathrm i\xi x}-e^{2\pi\mathrm i\eta x}|
 \le2\pi M|\xi-\eta|.
\]
这证明参数到字符的映射连续。反方向令 \(h=\xi-\eta\)。若
\[
 \sup_{|x|\le1}|e^{2\pi\mathrm ihx}-1|<\varepsilon<1,
\]
则必有 \(|h|<1/2\)：否则取 \(x=1/(2h)\)，左侧上确界至少为 \(2\)。在 \(|h|<1/2\) 时，再取 \(x=1\) 并用 \(\sin(\pi s)\ge2s\)（\(0\le s\le1/2\)），可得
\[
 4|h|\le2|\sin(\pi h)|=|e^{2\pi\mathrm ih}-1|<\varepsilon.
\]
因此紧集上的一致接近也迫使参数接近，逆映射连续。

\(\mathbb Z\) 中的紧集恰为有限集。若 \(z,w\in\mathbb S^1\)，整数幂的等比和公式给出
\[
 |z^n-w^n|\le |n|\,|z-w|\qquad(n\in\mathbb Z).
\]
所以 \(z\) 接近 \(w\) 时，相应字符在每个有限集合上一致接近；逆映射则是连续的求值 \(\chi\mapsto\chi(1)\)。这证明 \(\widehat{\mathbb Z\kern0pt}\simeq\mathbb T\)。

在 \(\mathbb T\) 上，若 \(n\ne m\)，取 \(t=1/(2(n-m))\) 的等价类，得到
\[
 \sup_{t\in\mathbb T}|e^{2\pi\mathrm int}-e^{2\pi\mathrm imt}|=2.
\]
整个 \(\mathbb T\) 是紧集，因此 \(U(\chi_n;\mathbb T,1)=\{\chi_n\}\)。每个字符都是孤立点，这正是 \(\mathbb Z\) 的离散拓扑。有限循环群的字符集是有限 Hausdorff 空间，因而同样离散。

最后考虑有限乘积。限制到某个坐标轴是连续映射，因为该坐标轴上紧集的像仍紧。反方向，若 \(K\subset\prod_jG_j\) 紧，则各坐标投影 \(K_j\) 紧，而且
\[
 \sup_{x\in K}\left|\prod_j\chi_j(x_j)-\prod_j\psi_j(x_j)\right|
 \le\sum_{j=1}^d\sup_{x_j\in K_j}|\chi_j(x_j)-\psi_j(x_j)|.
\]
逐项相减得到的这个估计证明逆方向的连续性，完成乘积拓扑的识别。
\end{proof}

\cref{b2:tab:I3-frequency-types}概括了这些模型中“位置”与“频率”的关系。周期位置对应整数频率；离散位置对应一个紧的周期频率变量；实直线则在两个方向上都保留连续参数。
\begin{table}[htbp]
\centering
\caption{基本交换群与它们的频率空间}
\label{b2:tab:I3-frequency-types}
\begin{tabular}{lll}
\toprule
位置群 \(G\) & 对偶群 \(\widehat G\) & 基本字符 \\
\midrule
\(\mathbb Z/N\mathbb Z\) & \(\mathbb Z/N\mathbb Z\)（有限） & \(e^{2\pi\mathrm ikj/N}\) \\
\(\mathbb Z^d\) & \(\mathbb T^d\)（紧） & \(e^{2\pi\mathrm in\cdot\theta}\) \\
\(\mathbb T^d\) & \(\mathbb Z^d\)（离散） & \(e^{2\pi\mathrm in\cdot t}\) \\
\(\mathbb R^d\) & \(\mathbb R^d\) & \(e^{2\pi\mathrm ix\cdot\xi}\) \\
\bottomrule
\end{tabular}
\end{table}

对偶空间中拓扑的差别具有实际计算意义。在 \(\mathbb T\) 上，两个不同整数频率在整个位置空间内总能积累到完全相反的相位；在 \(\mathbb R\) 上，相近频率只有经过足够长的位置区间才会积累出同样的差异。紧开拓扑恰好保留这种有限观测范围内的区别。

\subsection{Fourier 变换与卷积公式}
\label{b2:subsec:I20303}

\begin{definition}[局部紧交换群上的 Fourier 变换]
\label{b2:def:I3-fourier}
设 \(G\) 为第二可数局部紧 Hausdorff 交换群，\(\mu\) 是其固定的非零 Haar 测度。对 \(f\in L^1(G,\mu)\)，定义
\[
 \widehat f(\chi)=\int_G f(x)\overline{\chi(x)}\,\mathrm d\mu(x),
 \qquad \chi\in\widehat G.
\]
这个复值函数称为 \(f\) 的 \term[qunshangFourierbianhuan]{群上的 Fourier 变换}[Fourier transform on a group]。
\end{definition}

字符的模恒为 \(1\)，所以定义中的积分绝对收敛，并且不随 \(f\) 的几乎处处代表元改变。这里的频率变量是函数 \(\chi\)，在 \(\mathbb R^d\) 的具体识别下才成为通常的向量 \(\xi\)。测度的归一化也进入定义：把 \(\mu\) 改成 \(c\mu\) 会使 \(\widehat f\) 改成 \(c\widehat f\)。反演公式中的对偶测度将在下一节与此配对。

\begin{proposition}[有界性与连续性]
\label{b2:prop:I3-fourier-continuity}
设 \(G\) 为第二可数局部紧 Hausdorff 交换群，\(\mu\) 为非零 Haar 测度。若 \(f\in L^1(G,\mu)\)，则 \(\widehat f\) 是对偶群紧开拓扑上的有界连续函数，且
\[
 \|\widehat f\|_\infty\le\|f\|_1.
\]
因此 \(f\mapsto\widehat f\) 是从 \(L^1(G,\mu)\) 到 \(C_b(\widehat G)\) 的有界线性映射。
\end{proposition}
\begin{proof}
绝对值估计直接给出有界性。为证明在固定 \(\chi_0\) 处连续，先选紧集 \(K\subset G\)，使
\(\int_{G\setminus K}|f|\,\mathrm d\mu\) 足够小。这种选择来自 Haar 测度的正则性；也可以用第二可数局部紧空间的递增紧穷竭，再对 \(|f|\) 使用单调收敛。对任意字符 \(\chi\)，有
\[
 |\widehat f(\chi)-\widehat f(\chi_0)|
 \le\|f\|_1\sup_{x\in K}|\chi(x)-\chi_0(x)|
      +2\int_{G\setminus K}|f|\,\mathrm d\mu.
\]
给定误差 \(\varepsilon>0\)，先令右边尾项小于 \(\varepsilon/2\)，再在紧开邻域中令第一项小于 \(\varepsilon/2\)。这给出了连续性的邻域证明；当 \(f=0\) 时结论直接成立。
\end{proof}

定义中的共轭固定了全部相位符号。把平移和乘以字符分别带入积分，可以逐项确定频率端发生了什么。

\begin{proposition}[平移、调制、卷积与对合]
\label{b2:prop:I3-fourier-calculus}
设 \(G\) 为第二可数局部紧 Hausdorff 交换群，\(\mu\) 为非零 Haar 测度，\(f,g\in L^1(G,\mu)\)，\(a\in G\)，\(\eta,\chi\in\widehat G\)。定义
\[
 L_af(x)=f(x-a),\qquad M_\eta f(x)=\eta(x)f(x),\qquad
 f^*(x)=\overline{f(-x)}.
\]
则这些函数都属于 \(L^1(G,\mu)\)，且
\begin{align}
 \widehat{L_af}(\chi)&=\overline{\chi(a)}\,\widehat f(\chi),
    \label{b2:eq:I3-translation}\\
 \widehat{M_\eta f}(\chi)&=\widehat f(\chi\eta^{-1}),
    \label{b2:eq:I3-modulation}\\
 \widehat{f*g}(\chi)&=\widehat f(\chi)\widehat g(\chi),
    \label{b2:eq:I3-convolution}\\
 \widehat{f^*}(\chi)&=\overline{\widehat f(\chi)}.
    \label{b2:eq:I3-involution}
\end{align}
其中卷积使用固定的 Haar 测度，即
\((f*g)(x)=\int_G f(y)g(x-y)\,\mathrm d\mu(y)\)。
\end{proposition}
\begin{proof}
平移保持 Haar 测度，乘以字符保持模，交换群上的取逆 \(x\mapsto-x\) 也保持 Haar 测度，故所列三个函数与 \(f\) 有相同的 \(L^1\) 范数。卷积属于 \(L^1\) 以及 \(\|f*g\|_1\le\|f\|_1\|g\|_1\) 已在\cref{b2:thm:I2-convolution}证明。

在平移积分中令 \(x=a+y\)，得到
\[
 \int_G f(x-a)\overline{\chi(x)}\,\mathrm d\mu(x)
 =\overline{\chi(a)}\int_G f(y)\overline{\chi(y)}\,\mathrm d\mu(y).
\]
调制公式来自点态等式
\(\eta(x)\overline{\chi(x)}=\overline{(\chi\eta^{-1})(x)}\)。

卷积公式所需的换序可由以下绝对可积性核实：
\[
 \int_G\int_G|f(y)|\,|g(x-y)|\,\mathrm d\mu(x)\,\mathrm d\mu(y)
 =\|f\|_1\|g\|_1<\infty.
\]
因此 Fubini 定理适用，再令 \(z=x-y\)，便得
\begin{align*}
 \widehat{f*g}(\chi)
 &=\int_G\int_G f(y)g(x-y)\overline{\chi(x)}
       \,\mathrm d\mu(x)\,\mathrm d\mu(y)\\
 &=\int_G\int_G f(y)g(z)\overline{\chi(y)}\,
       \overline{\chi(z)}\,\mathrm d\mu(z)\,\mathrm d\mu(y)\\
 &=\widehat f(\chi)\widehat g(\chi).
\end{align*}
最后，在对合积分中令 \(x=-y\)，利用 \(\overline{\chi(-y)}=\chi(y)\)，得到
\[
 \widehat{f^*}(\chi)
 =\int_G\overline{f(y)}\chi(y)\,\mathrm d\mu(y)
 =\overline{\widehat f(\chi)}.
\]
\end{proof}

例如，若 \(g\in L^1(G)\)，卷积算子 \(T_gf=g*f\) 在每个频率上只是乘以 \(\widehat g(\chi)\)。这使卷积方程 \(g*f=h\) 产生代数关系
\(\widehat g\,\widehat f=\widehat h\)。当 \(\widehat g\) 有零点时，哪些 \(h\) 能被解出立即受到约束；即使没有零点，频率上除法所得的函数是否对应所需空间中的 \(f\)，仍要结合反演与范数估计判断。下一节的反演和 Plancherel 公式将给出这种判断的基本工具。
